\chapter[\emj{🗝️}~Patrones con Hash Map (Frecuencias, Prefix-Sum, Agrupación)]{\emj{🗝️}~Patrones con Hash Map}

\begin{dsaquees}

El Hash Map es la navaja suiza 🔪 de las entrevistas: cuando algo tarda
$O(n^2)$ "buscando de nuevo", casi siempre un Hash Map lo baja a $O(n)$
recordando lo que ya viste.

Tres patrones que aparecen una y otra vez:
\begin{itemize}
\item Contar frecuencias: cuántas veces aparece cada cosa.
\item Complemento (Two Sum): guardar lo que YA viste para encontrar su
      pareja al instante.
\item Prefix-Sum + Hash: contar sub-arreglos cuya suma es un objetivo,
      recordando las sumas acumuladas.
\end{itemize}

\end{dsaquees}

\begin{dsaotraforma}

El Hash Map convierte "¿ya vi esto antes?" en una consulta $O(1)$. Los
patrones más rentables:

\begin{itemize}
\item \textbf{Frequency Count}: un mapa \verb|valor -> conteo|. Base de
      anagramas, mayoría, top-K, duplicados.
\item \textbf{Two Sum / complemento}: por cada $x$ busca si $target-x$ ya
      está en el mapa. Una sola pasada, $O(n)$.
\item \textbf{Group By (agrupación)}: mapa \verb|clave_canónica -> lista|.
      Para anagramas la clave es la palabra ordenada o su vector de
      conteo de letras.
\item \textbf{Prefix-Sum + Hash}: para "subarreglos que suman K", guarda
      cuántas veces apareció cada suma acumulada; si \verb|suma - K| ya
      se vio, hay subarreglos que suman K.
\end{itemize}

\end{dsaotraforma}

\begin{dsadiagrama}
\begin{verbatim}
TWO SUM (target = 9)          SUBARRAY SUM = K (K = 3)

nums = [2, 7, 11, 15]         nums = [1, 2, 1, 3], K = 3
visto = {}                    prefix acumulada: 1, 3, 4, 7
x=2 -> falta 7? no. guarda 2  mapa {suma_prefija -> veces}, inicia {0:1}
x=7 -> falta 2? SÍ (indice 0) suma=1 -> busca 1-3=-2? no. mapa{0:1,1:1}
       => par (0,1)           suma=3 -> busca 3-3=0? SÍ (1 vez) -> +1
                              suma=4 -> busca 4-3=1? SÍ (1 vez) -> +1
                              suma=7 -> busca 7-3=4? SÍ (1 vez) -> +1
                              total subarreglos con suma 3: 3
\end{verbatim}
\end{dsadiagrama}

\begin{dsacomofunciona}
\begin{verbatim}
TWO SUM   nums=[2,7,11,15]  target=9   (mapa valor→índice)

 i │ nums[i] │ falta=9-x │ ¿falta en mapa? │ mapa después
 ──┼─────────┼───────────┼─────────────────┼──────────────────
 0 │   2     │    7      │ no              │ {2:0}
 1 │   7     │    2      │ SÍ (índice 0)   │ → respuesta (0,1) OK

SUBARRAY SUM = K   nums=[1,2,1,3]  K=3   mapa{sumaPrefija:veces}

 x │ suma │ busca suma-K │ ¿está? +cuenta │ mapa (inicia {0:1})
 ──┼──────┼──────────────┼────────────────┼─────────────────────
 1 │  1   │     -2       │ no             │ {0:1, 1:1}
 2 │  3   │      0       │ SÍ ×1  total=1 │ {0:1, 1:1, 3:1}
 1 │  4   │      1       │ SÍ ×1  total=2 │ {0:1, 1:1, 3:1, 4:1}
 3 │  7   │      4       │ SÍ ×1  total=3 │ {..., 7:1}

Total subarreglos con suma 3 = 3   ([1,2] [3] [2,1] desde prefijos)
Clave: el mapa "recuerda" sumas ya vistas → cada consulta O(1).
\end{verbatim}
\end{dsacomofunciona}

\begin{lstlisting}
TWO SUM
- Un mapa valor->índice. Por cada x, si target-x ya está, encontraste
  el par. Si no, guarda x. Una sola pasada.

SUBARRAY SUM = K
- Suma acumulada 'suma'. Un mapa 'cuantas veces salió cada suma'.
- Si (suma - K) está en el mapa, esos son subarreglos que suman K.
- Inicializa el mapa con {0: 1} (el prefijo vacío).

📝 PSEUDOCÓDIGO (subarray sum = K)
──────────────────────────────────────────────────────────────

función subarraysConSuma(nums, K):
    conteo = {0: 1}; suma = 0; total = 0
    PARA cada x en nums:
        suma += x
        SI (suma - K) en conteo: total += conteo[suma - K]
        conteo[suma] += 1
    return total

💻 CÓDIGO C++ (Two Sum)
──────────────────────────────────────────────────────────────

#include <vector>
#include <unordered_map>
using namespace std;

vector<int> twoSum(vector<int>& nums, int target) {
    unordered_map<int,int> visto;              // valor -> índice
    for (int i = 0; i < (int)nums.size(); i++) {
        int falta = target - nums[i];
        if (visto.count(falta)) return {visto[falta], i};
        visto[nums[i]] = i;
    }
    return {};
}

💻 CÓDIGO C++ (subarray sum = K)
──────────────────────────────────────────────────────────────

int subarraySum(vector<int>& nums, int k) {
    unordered_map<int,int> conteo{{0, 1}};     // prefijo vacío
    int suma = 0, total = 0;
    for (int x : nums) {
        suma += x;
        if (conteo.count(suma - k)) total += conteo[suma - k];
        conteo[suma]++;
    }
    return total;
}

💻 CÓDIGO C++ (agrupar anagramas)
──────────────────────────────────────────────────────────────

#include <string>
#include <algorithm>
vector<vector<string>> groupAnagrams(vector<string>& strs) {
    unordered_map<string, vector<string>> grupos;
    for (auto& s : strs) {
        string clave = s;
        sort(clave.begin(), clave.end());      // clave canónica
        grupos[clave].push_back(s);
    }
    vector<vector<string>> res;
    for (auto& [k, v] : grupos) res.push_back(v);
    return res;
}

⏱️ COMPLEJIDAD (Big-O)
──────────────────────────────────────────────────────────────

  Patrón            │ Tiempo       │ Espacio
  ──────────────────┼──────────────┼──────────
  Frequency / TwoSum│ O(n)         │ O(n)
  Subarray sum = K  │ O(n)         │ O(n)
  Group anagrams    │ O(n·k log k) │ O(n·k)

* k = longitud de la palabra en group anagrams.
\end{lstlisting}

\begin{dsaentrevistas}

✅ Regla mental: si te descubres haciendo un bucle DENTRO de otro para
   "volver a buscar", pregúntate "¿un Hash Map recuerda esto en O(1)?".
   Suele bajar O(n²) a O(n).

💡 Prefix-Sum + Hash es el patrón estrella oculto: "Subarray Sum Equals
   K" (LeetCode 560), "Continuous Subarray Sum", "Contiguous Array".
   Inicializar el mapa con \verb|{0: 1}| es el detalle que a todos se les
   olvida.

⚠️ Group anagrams por conteo de letras: en vez de ordenar (k log k),
   usa un vector de 26 conteos como clave -> O(k). Mejora que impresiona.

🔑 Hash Set vs Hash Map: si solo importa "¿existe?" usa un Set; si
   necesitas asociar un dato (índice, conteo) usa un Map.

\end{dsaentrevistas}

\begin{dsaresumen}

• Hash Map convierte "¿ya lo vi?" en O(1): rompe bucles O(n²).
• Frequency count: base de anagramas, mayoría, duplicados, top-K.
• Two Sum: guarda visto, busca el complemento target−x.
• Prefix-Sum + Hash: subarreglos con suma K; inicia el mapa con {0:1}.
• Group by: mapa de clave canónica -> lista.
• Set si solo importa existencia; Map si asocias datos.
════════════════════════════════════════════════════════════════

\end{dsaresumen}
